% Generado por gen_tabla_heterogeneidad.py. No editar a mano.

\begin{table}[H]
    \centering
    \begingroup
    \footnotesize
    \setlength{\tabcolsep}{4pt}
    \renewcommand{\texttt}[1]{{\ttfamily #1}}
    \ttabbox[\FBwidth]{
        \caption{Heterogeneidad de la anotación humana, en porcentaje de piezas marcadas como \textit{Sí}. Las tres primeras columnas comparan a los dos equipos y las tres últimas recorren a las anotadoras una a una. Todas las cifras se restringen a las 7 personas anotadoras que codificaron al menos 50 piezas (1.300 de las 1.313). Ningún modelo interviene en estas cifras.}
        \label{tab:iris-heterogeneidad}
    }{
        \begin{tabular*}{\textwidth}{@{\extracolsep{\fill}}l cc c cc >{\columncolor{gray!15}}c@{}}
            \toprule
            & \multicolumn{3}{c}{\textbf{Entre equipos}} & \multicolumn{3}{c}{\textbf{Entre anotadoras}} \\
            \cmidrule(lr){2-4} \cmidrule(lr){5-7}
            \textbf{Variable} & Indexa & UCM3 & \textbf{Brecha} & Mín. & Máx. & \textbf{Rango} \\
            \midrule
            V25 \texttt{lenguaje\_sexista} & 75,0 & 85,5 & \textbf{10,5} & 66,8 & 98,6 & \textbf{31,8} \\
            V26 \texttt{masc\_generico} & 71,8 & 88,1 & \textbf{16,3} & 64,9 & 97,5 & \textbf{32,6} \\
            V30 \texttt{sexismo\_discurso} & 11,4 & 65,2 & \textbf{53,8} & 3,9 & 97,8 & \textbf{93,9} \\
            V33 \texttt{asimetria\_mujer\_hombre} & 3,4 & 8,2 & \textbf{4,9} & 0,9 & 27,1 & \textbf{26,2} \\
            V35 \texttt{denominacion\_sexualizada} & 6,3 & 12,4 & \textbf{6,1} & 0,0 & 21,5 & \textbf{21,5} \\
            \bottomrule
        \end{tabular*}
    }
    \endgroup
\end{table}
